% !PRINTED 383
\weberpage{383}
%% [running head, not body text -- commented out 2026-08-07] Page furniture (running
%% head, folio, signature) is excluded from this edition by standing policy.
%% §\,104. Characters and genera of quadratic forms. \hfill 383

If the class $A$ is compounded from the class $A'$, $A''$, then
symbolically
\[
A = A'A''.
\]

Let $m'$, $m''$ be representable by the class $A'$, $A''$, and let us
take $m'$, $m''$ relatively prime; then $m = m'm''$ is representable by $A$
(§ 103), and from this there results for the characters the formula:

\weq{13}{\chi(\delta, A'A'') = \chi(\delta, A')\chi(\delta, A'').}

By 1., the symbol $\chi(\delta, A)$ is determined by any number $m$
coprime to $\delta$ and properly representable by $A$,
and we therefore define the characters of these numbers $m$ by

\weq{14}{\chi(\delta, m) = \chi(\delta, A).}

It is, however, expedient to free oneself also from the assumption that $m$ be relatively prime to $\delta$, and this
can be done in the following way. The assumption that
$m$ be relatively prime to $Q$ is, however, to be retained.

If $m$ is properly representable by the class $A$, then we
can find in $A$ a representative as follows:

\weq{15}{\varphi = (m, B, C)}

We decompose $m$ into two factors

\weq{16}{m = nn',}

so that $n$ and $n'$ are coprime and $n$ is relatively prime to $\delta$, $n'$
relatively prime to $\delta'$.

This is always possible, and mostly in several ways, since by
assumption $m$ is relatively prime to $Q$, and therefore $m$, $\delta$, $\delta'$ have no
common divisor. One may, for example, take into $n$ the primes
dividing $\delta'$ in powers as high as those in
which they are contained in $m$, and besides these the prime powers
that are coprime to $\delta$.

Then the form $\varphi$ is compounded from the two forms
\[
(n, B, Cn') \qquad (n', B, Cn)
\]
whose characters, by (12) and 4., are determined by
\[
(\delta, n), \quad (\delta'\ n') % [sic] German source prints no comma in the second symbol.
\]
and hence, by (13):

\weq{17}{\chi(\delta, A) = (\delta, n)(\delta', n').}

In this, however, the requirement that $m$ be properly
representable by $A$ can also be given up; for a
common divisor of $x, y$ (relatively prime to $Q$) gives a
%CONT>>
